\section{习题 3.4：二次曲线的仿射特征}
\chr{24}{163—166}{169—172}{305—307}
\K{本节的工作公式：中心 $h$ 解 $Ah=-b$；渐近方向 $u=(u_x,u_y)$ 解 $u^\top Au=0$；关于中心的坐标 $X=x-h_x,Y=y-h_y$ 会消去所有一次项；过曲线上正则点 $P$ 的切线为 $\nabla F(P)\cdot(X-P)=0$。}
\begin{center}
\begin{tikzpicture}[scale=.78]
\draw[->](-3,0)--(3.4,0)node[right]{$X$};\draw[->](0,-2.25)--(0,2.25)node[above]{$Y$};
\draw[dashed,accent,thick](-2,-2)--(2,2) node[above right]{$L_1$};
\draw[dashed,accent,thick](-2,2)--(2,-2) node[below right]{$L_2$};
\draw[very thick,mint,domain=-1.2:1.2,samples=90]plot({sqrt(1+\x*\x)},\x);
\draw[very thick,mint,domain=-1.2:1.2,samples=90]plot({-sqrt(1+\x*\x)},\x);
\fill (0,0)circle(1.5pt) node[below left]{$H$};
\node[align=left]at(4.2,1.5){双曲线的渐近线：\\过中心 $H$，方向满足 $u^\top Au=0$};
\end{tikzpicture}
\end{center}

\Q{1}{十条二次曲线：逐项求中心、渐近方向与渐近线}
\hint{对每一式先算 $A$ 与 $b$；$\det A\ne0$ 时解中心；$\det A<0$ 时令 $\Phi(X,Y)=0$ 分解成渐近线。$\det A=0$ 时再判抛物线还是平行线，不能强行使用 $A^{-1}$。}
\subpart{(1) $x^2-2xy-4y^2+6x+2y+3=0$} 
$A=\begin{psmallmatrix}1&-1\\-1&-4\end{psmallmatrix}$、$b=(3,1)^\top$。解 $Ah=-b$ 得 $H=(-11/5,4/5)$。移到中心后
\[X^2-2XY-4Y^2-\tfrac{14}{5}=0.\]
解 $\Phi(u)=u_x^2-2u_xu_y-4u_y^2=0$，方向向量可取 $(1+\sqrt5,1)$、$(1-\sqrt5,1)$，过 $H$ 的渐近线为
\[
\begin{cases}5x-5(1+\sqrt5)y+15+4\sqrt5=0,\\5x-5(1-\sqrt5)y+15-4\sqrt5=0.\end{cases}
\]
\subpart{(2) $3x^2+4xy+4y^2+6x+2y-12=0$} 
$A=\begin{psmallmatrix}3&2\\2&4\end{psmallmatrix}$、$b=(3,1)^\top$，$\det A=8>0$ 且正定。$H=(-5/4,3/8)$，$F(H)=-123/8<0$，因此实椭圆，\textbf{没有实渐近方向和渐近线}。
\subpart{(3) $x^2-2xy+y^2+x+2y+4=0$} 
$F=(x-y)^2+x+2y+4$，$A$ 秩 1，一次项沿零方向 $(1,1)$ 非零，无中心。只有渐近方向 $(1,1)$；沿 $(1,1)$ 增加参数会使一次项增加，故开口朝向 $(-1,-1)$。该曲线是抛物线，\textbf{不存在通常双曲线意义的有限渐近线}。
\subpart{(4) $x^2+2xy+y^2+2x+2y-3=0$}
直接分解：$(x+y-1)(x+y+3)=0$，两条平行直线。其共同方向为 $(1,-1)$；中心不唯一（中线 $x+y+1=0$ 上各点都为中心），两条直线自身可视为退化曲线的两支。
\subpart{(5) $5x^2-12xy+4y^2-4x-8y+4=0$}
$H=(-2,-2)$，中心化 $F=5X^2-12XY+4Y^2+16$；
\[5X^2-12XY+4Y^2=(5X-2Y)(X-2Y).
\]
方向取 $(2,5)$ 与 $(2,1)$，渐近线 $5x-2y+6=0$、$x-2y-2=0$。
\subpart{(6) $4xy-5y^2-12x-6y-11=0$}
$H=(9,3)$，中心化 $F=4XY-5Y^2-74$；渐近方向 $(1,0)$、$(5,4)$，渐近线 $y-3=0$、$4x-5y-21=0$。
\subpart{(7) $5x^2+4xy+y^2-6x+4y-6=0$}
$H=(7,-16)$，$A=\begin{psmallmatrix}5&2\\2&1\end{psmallmatrix}$ 正定、中心处 $F(H)=-59$；实椭圆，无实渐近线。
\subpart{(8) $3x^2+4xy+5y^2-2x+6y-1=0$}
$H=(1,-1)$，$A=\begin{psmallmatrix}3&2\\2&5\end{psmallmatrix}$ 正定、$F(H)=-5$；实椭圆，无实渐近线。
\subpart{(9) $5x^2+8xy+3y^2+8x+6y=0$}
$H=(0,-1)$，中心化 $F=5X^2+8XY+3Y^2-3$；
\[(X+Y)(5X+3Y)=0.
\]
渐近方向 $(1,-1)$、$(3,-5)$；渐近线 $x+y+1=0$、$5x+3y+3=0$。
\subpart{(10) $8x^2+8xy+2y^2-6x-4y+5=0$}
$F=2(2x+y)^2-6x-4y+5$。$A$ 秩 1，没有中心；其零特征方向为 $(1,-2)$。沿 $(-1,2)$ 的一次项使 $F$ 降低，所以抛物线开口朝 $(-1,2)$。
\ans{十项的中心、渐近方向与渐近线（如果存在）均已逐式得到，并与原书列出的数值和方向核对。}\readtip{\textbf{最容易漏：}第 (4) 项不是一条抛物线而是两条平行线；第 (3)、(10) 项的“渐近方向”不能误当成一条有限渐近线。}

\Q{2}{参数 $s,t$ 对中心和渐近方向的影响}
\textbf{曲线：}$x^2+6xy+sy^2+3x+ty-4=0$。\hint{中心方程 $x+3y=-3/2,\ 3x+sy=-t/2$；渐近方向解 $u_x^2+6u_xu_y+su_y^2=0$。}
\subpart{中心} 系数行列式为 $s-9$。$s\ne9$ 时有唯一中心；$s=9,t\ne9$ 时两方程矛盾，无中心；$s=t=9$ 时两式等价，有一条\textbf{中心线} $x+3y=-3/2$。此时曲线恰为两平行线 $(x+3y+4)(x+3y-1)=0$。
\subpart{渐近方向} $\Phi(u)=(u_x+3u_y)^2+(s-9)u_y^2$。若 $s<9$，有两条实渐近方向 $u_x=(-3\pm\sqrt{9-s})u_y$；若 $s=9$，方向合并为一条 $(-3,1)$；若 $s>9$，没有实零方向。$t$ 只在一次项中，因此不影响这些方向。
\ans{中心存在条件与渐近方向的三种情况均完成。}

\Q{3}{过三点且中心已知的二次曲线}
\textbf{题意：}过 $(2,3),(4,2),(1,5)$ 且 $H=(0,1)$ 是中心。\hint{移中心：令 $X=x,Y=y-1$，无一次项，设 $aX^2+2bXY+cY^2+d=0$。}
三点变成 $(2,2),(4,1),(1,4)$，代入
\[
4a+8b+4c+d=0,\quad16a+8b+c+d=0,\quad a+8b+16c+d=0.
\]
由第一式减第二式得 $c=4a$，第一减第三得 $a=4c$，从而 $a=c=0,d=-8b$。非零比例取 $2b=1$，则
\[\boxed{x(y-1)-4=0\quad\Longleftrightarrow\quad2xy-2x-8=0.}\]
\ans{方程与原书答案一致；它实际是一条以 $(0,1)$ 为中心的双曲线。}

\Q{4}{中心点值与第三、第二不变量的关系}
\textbf{题意：}若 $h$ 是曲线 $F=0$ 的中心，证明 $I_3=I_2F(h)$。\hint{以中心为新原点做平移，对二次型增广矩阵做合同变换。}
$Ah=-b$ 意味着 $F(h+X)=X^\top AX+F(h)$。平移对应齐次矩阵 $T=\begin{psmallmatrix}I&h\\0&1\end{psmallmatrix}$，满足 $\det T=1$，于是
\[
T^\top H T=\begin{pmatrix}A&0\\0&F(h)\end{pmatrix}
\Longrightarrow \det H=\det A\,F(h).
\]
\ans{$I_3=I_2F(h)$；此证明甚至不需要假定 $A$ 一定可逆。}

\Q{5}{共轭直径给定时的二次曲线}
\textbf{题意：}曲线过 $(-2,-1),(0,-2)$，$x+y+1=0,x-y+1=0$ 为一对共轭直径。\hint{中心是两直径交点 $H=(-1,0)$，在中心坐标 $X=x+1,Y=y$ 下设 $aX^2+2bXY+cY^2+d=0$。}
两直径方向 $(1,-1),(1,1)$ 互相共轭，故 $B_A((1,-1),(1,1))=a-c=0$，即 $c=a$。带入两个已知点：
\[
2a+2b+d=0,\qquad5a-4b+d=0
\Longrightarrow b=a/2,\ d=-3a.
\]
令 $a=1$，得到 $X^2+XY+Y^2-3=0$；还原
\[\boxed{x^2+y^2+xy+2x+y-2=0.}\]
\ans{与书后结果一致；与 3.1 第 10 题所得方程相同，但两题给的是不同几何数据。}

\Q{6}{已知两条渐近线和一个点确定双曲线}
\textbf{条件：}$y-2=0$、$3x+4y+5=0$ 为渐近线，且过 $(0,1/2)$。\hint{与两条渐近线对应的双曲线方程可写成其线性形式乘积加常数。}
设 $F=(y-2)(3x+4y+5)+k$。带入 $(0,1/2)$：$(-3/2)\cdot7+k=0$，得 $k=21/2$。展开后乘 2：
\[\boxed{6xy+8y^2-12x-6y+1=0.}\]
\ans{与书后答案一致；原题给出的是渐近线而不是对称轴。}

\Q{7}{原点为二次曲线中心的系数条件}
\hint{对中心 $h=0$ 使用 $Ah=-b$。}
$A0=-b$ 强制 $b=0$，即 $b_1=b_2=0$；也可比较 $F(x,y)$ 和 $F(-x,-y)$，二者相等意味着一次项全为零。
\ans{\textbf{$b_1=b_2=0$}，其他二次项和常数项不因此受限。}

\Q{8}{渐近线是坐标轴的双曲线特征}
\hint{渐近线相交于中心，故 $b_1=b_2=0$；两坐标轴方向 $e_x,e_y$ 都要使 $\Phi$ 为零。}
$\Phi(1,0)=a_{11}=0$，$\Phi(0,1)=a_{22}=0$，剩下
\[F=2a_{12}xy+c.
\]
这是以两坐标轴为渐近线的非退化双曲线，当且仅当 $a_{12}\ne0,c\ne0$；若 $c=0$ 就变成两条相交直线而非真双曲线。
\ans{$a_{11}=a_{22}=b_1=b_2=0,\ a_{12}c\ne0$。}

\Q{9}{一条双曲线经过原点的直径与它的共轭直径}
\textbf{条件：}$x^2-2xy-4y^2+6x+2y+3=0$（同第 1(1) 项）。\hint{任何直径都经过中心 $H=(-11/5,4/5)$；过原点的直径唯一为 $OH$；共轭条件为 $u^\top Av=0$。}
$OH$ 的方向取 $u=(-11,4)$，过原点直径为
\[\boxed{4x+11y=0.}
\]
取 $A=\begin{psmallmatrix}1&-1\\-1&-4\end{psmallmatrix}$，得 $Au=(-15,-5)$，共轭向量须满足 $-15v_x-5v_y=0$，可取 $v=(1,-3)$。共轭直径经 $H$，其方程
\[\boxed{15x+5y+29=0.}\]
\ans{与书后答案一致。}

\Q{10}{弦中点对应的直径：一般证明及三组算例}
\subpart{(1) 一般证明} 设弦端点 $P\pm v$ 均在二次曲线 $F=0$ 上。二次函数恒满足
\[F(P+v)-F(P-v)=2\nabla F(P)\cdot v=0.
\]
故弦所在直线的方向 $v$ 与 $\nabla F(P)$ 垂直，所以弦的直线方程为
\[\boxed{(x-x_0)F_1(x_0,y_0)+(y-y_0)F_2(x_0,y_0)=0},
\]
这里课本用 $F_1=\tfrac12\partial_xF,\;F_2=\tfrac12\partial_yF$；式子乘 2 不影响直线。
\subpart{(2) (1)} $F=5x^2+8xy+5y^2-18x-18y+9$，$P=(1,0)$；
\[F_1(P)=5+4\cdot0-9=-4,\quad F_2(P)=4+5\cdot0-9=-5
\Rightarrow\boxed{4x+5y-4=0}.\]
\subpart{(2)} $F=x^2+2xy+y^2-8x+4$，$P=(2,1)$；
\[F_1(P)=2+1-4=-1,\quad F_2(P)=2+1=3
\Rightarrow\boxed{x-3y+1=0}.\]
\subpart{(3)} $F=6xy+8y^2-12x-26y+11$，$P=(3,1)$；
\[F_1(P)=3-6=-3,\quad F_2(P)=9+8-13=4
\Rightarrow\boxed{3x-4y-5=0}.\]
\caution{第 (3) 项的最后简化不能直接用 $(-3,4)\cdot(x-3,y-1)=0$ 得出 $x-3=0$。实际从这个点出发的直径应为 $3x-4y-5=0$；但原书答案第 (3) 项写为 $x-3=0$。\textbf{请核查原题系数：}若原题 $F$ 为 $6xy+8y^2-12x-26y+11$，则上面计算严格推出 $3x-4y-5=0$，而不是书后列值。}
\ans{(1) $4x+5y-4=0$；(2) $x-3y+1=0$；(3) 按题干系数应为 $3x-4y-5=0$（书后印为 $x-3=0$，已标记复核）。}

\Q{11}{直角坐标轴是共轭直径时的标准形式}
\hint{既然两坐标轴都是直径，中心在原点，故一次项为零；轴向量 $e_1,e_2$ 共轭又要求 $e_1^\top A e_2=a_{12}=0$。}
因此 $b_1=b_2=a_{12}=0$，曲线就是
\[\boxed{a_{11}x^2+a_{22}y^2+c=0}.\]
\ans{得证；椭圆、双曲线等标准形都是这个结论的特例。}

\Q{12}{渐近方向与平行渐近线的关系}
\hint{以双曲线中心为原点，渐近线由二次齐次部分 $\Phi(X,Y)=0$ 给出。}
任何非零方向 $u$ 若满足 $\Phi(u)=0$，沿该方向的射线 $X=tu$ 是二次曲线的主导零方向，其对应的一条过中心渐近线就是 $X=tu$。\textbf{方向向量 $u$ 不规定朝向}，所以对应的渐近线要用通过中心的直线而非半射线表达。
\ans{对非退化双曲线，任意渐近方向必对应平行于它的渐近线。}

\Q{13}{椭圆、双曲线与抛物线的“直径”性质}
\subpart{(1) 椭圆过中心的任一直线均为直径} 以中心为原点，椭圆为 $X^\top AX=1$，$A$ 正定。取任意非零直径方向 $u$，其 $A$-共轭方向可取非零 $v$ 满足 $u^\top Av=0$，以 $v$ 为弦方向的所有弦中点便落在过中心的 $u$ 直线上。因此任意过中心直线是某组平行弦的中点轨迹。
\subpart{(2) 双曲线同理} 双曲线 $X^\top AX=1$ 中 $A$ 不定。对非渐近方向 $u$，同样可取 $A$-共轭方向 $v$，得到过中心的直径；当涉及渐近方向，直径可视为相应的极限情形，需排除没有实际割线的方向。
\subpart{(3) 抛物线的直径} 取 $y^2=2px$。对方向 $v=(r,s)$ 且 $s\ne0$ 的弦，中点 $M=(x_0,y_0)$ 满足 $F_1(M)r+F_2(M)s=-pr+y_0s=0$，于是 $y_0=pr/s$，中点轨迹是平行于 $x$ 轴的直线。反过来任选 $y_0$，取 $r/s=y_0/p$ 可得到相应弦族。
\ans{椭圆和双曲线的直径通过中心；抛物线的每条直径都平行其轴方向，证明及例子如上。}

\Q{14}{椭圆一对共轭半径的平方和是常数}
\textbf{椭圆：}$x^2/a^2+y^2/b^2=1$。设共轭半径向量 $u,v$ 满足
\[u^\top Au=v^\top Av=1,\qquad u^\top Av=0,\quad A=\operatorname{diag}(1/a^2,1/b^2).
\]
令 $U=[u\ v]$，上述条件等价于 $U^\top AU=I$。由于 $U$ 可逆，有 $UU^\top=A^{-1}$。取迹便得
\[|u|^2+|v|^2=\operatorname{tr}(UU^\top)=\operatorname{tr}(A^{-1})=a^2+b^2.\]
\ans{\textbf{任意一对共轭半径的平方和恒为 $a^2+b^2$}；没有假设它们互相垂直。}

\Q{15}{六个给定曲线和给定点处的切线}
\hint{只要给定点在曲线上且 $\nabla F(P)\neq0$，取 $\nabla F(P)\cdot(X-P)=0$ 即可；以下每一小问均可直接代回验算。}
\begin{enumerate}
\item 椭圆 $x^2/a^2+y^2/b^2=1$，$P=(x_0,y_0)$：$\boxed{x_0x/a^2+y_0y/b^2=1}$。
\item 双曲线 $x^2/a^2-y^2/b^2=1$：$\boxed{x_0x/a^2-y_0y/b^2=1}$。
\item 抛物线 $y^2=2px$：$F_x=-2p,F_y=2y_0$，得 $\boxed{px-y_0y+px_0=0}$（等价于 $y_0y=p(x+x_0)$）。
\item $F=3x^2+4xy+5y^2-7x-8y-3$，$P=(2,1)$：$\nabla F(P)=(9,10)$，所以 $\boxed{9x+10y-28=0}$。
\item $F=5x^2+7xy+y^2-x+2y$，$P=(0,0)$：$\nabla F(P)=(-1,2)$，故 $\boxed{x-2y=0}$。
\item $F=5x^2-2xy+5y^2-8$，$P=(1,1)$：$\nabla F(P)=(8,8)$，得 $\boxed{x+y-2=0}$。
\end{enumerate}
\ans{六个切线结果均由点上切线公式直接验证，并与原书答案一致。}

\Q{16}{规定切线方向时求全部切点和切线（三组）}
\hint{令切线方向为 $u=(r,s)$。切点 $P$ 必须同时满足 $F(P)=0$ 和 $\nabla F(P)\cdot u=0$；后式是一次方程，消元通常给出两个点。}
\subpart{(1) $F=x^2+4xy+3y^2-5x-6y+3$，$u=(4,-1)$} 
\[
\nabla F\cdot u=(2x+4y-5)4-(4x+6y-6)=4x+10y-14.
\]
故 $2x+5y=7$。与 $F=0$ 联立得到 $P_1=(1,1)$、$P_2=(-4,3)$。切线都平行 $u$，法向 $(1,4)$，即
\[\boxed{x+4y-5=0,\qquad x+4y-8=0.}\]
\subpart{(2) $F=3x^2+4xy+5y^2-6x-8y+3$，$u=(1,0)$} 
切点处 $\nabla F\cdot u=6x+4y-6=0$；代入 $F=0$ 求得 $P_1=(1,0)$、$P_2=(3/11,12/11)$。相应切线均平行 $x$ 轴：
\[\boxed{y=0,\qquad 11y-12=0.}\]
\subpart{(3) $F=5x^2+8xy+y^2+2x+2y$，$u=(1,-1)$}
$\nabla F\cdot u=(10x+8y+2)-(8x+2y+2)=2x+6y=0$，即 $x+3y=0$。代入二次曲线得切点 $(0,0)$ 和 $(-6/11,2/11)$，切线法向 $(1,1)$：
\[\boxed{x+y=0,\qquad 11x+11y+4=0.}\]
\ans{三小题的\textbf{全部六个切点、六条切线}均列出，并逐个满足 $F(P)=0$ 和方向条件。}

\Q{17}{由轴方向、切点和另一个点确定一条抛物线}
\textbf{条件：}轴平行 $(1,-2)$，直线 $2x+13y-15=0$ 在 $(1,1)$ 处相切，且曲线经过 $(3,0)$。\hint{轴方向与二次型零特征向量一致；取 $\Phi=(2x+y)^2$，再求三个线性系数。}
设
\[F(x,y)=(2x+y)^2+Dx+Ey+G.
\]
曲线过 $(1,1)$、$(3,0)$，有 $9+D+E+G=0,\ 36+3D+G=0$。在 $(1,1)$ 的法向 $\nabla F=(12+D,6+E)$ 必须平行 $(2,13)$，所以 $13(12+D)=2(6+E)$。解得 $(D,E,G)=(-10,7,-6)$，即
\[\boxed{4x^2+4xy+y^2-10x+7y-6=0.}\]
\ans{代 $(1,1)$ 可验证切点条件，代 $(3,0)$ 可验证经过该点；与书后答案一致。}

\Q{18}{从外点引两条互相垂直切线的条件}
\textbf{题意：}二次曲线 $F$ 中课本记 $F_1=a_{11}x+a_{12}y+b_1$、$F_2=a_{12}x+a_{22}y+b_2$，证明外点 $P$ 可作两条正交实切线时满足 $F_1(P)^2+F_2(P)^2=I_1F(P)$。\hint{考察从 $P$ 沿方向 $u=(r,s)$ 发出的直线 $P+tu$ 与 $F=0$ 的交点重数。}
沿线代入得
\[F(P+tu)=\Phi(u)t^2+2\bigl(F_1(P)r+F_2(P)s\bigr)t+F(P).
\]
切线条件为 $[F_1(P)r+F_2(P)s]^2=F(P)\Phi(u)$。设两条切线方向是\textbf{正交单位向量} $u,v$，把各自等式相加：
\[
F_1(P)^2+F_2(P)^2=F(P)\bigl(\Phi(u)+\Phi(v)\bigr)=F(P)\operatorname{tr}A.
\]
故
\[\boxed{F_1(P)^2+F_2(P)^2=I_1F(P).}\]
更一般地，令 $T=gg^\top-F(P)A$，$g=(F_1,F_2)$；此等式等价于 $\operatorname{tr}T=0$。非零实对称无迹矩阵有两条互相垂直的零二次型方向，从而得到反向判别。但\textbf{若直线仅与曲线有无穷远交点}，仍须检查有有限实切点。
\ans{必要条件已严格证明，亦说明了反向应用的非退化检查步骤。}

\Q{19}{给定椭圆上正交切线交点的轨迹}
\textbf{曲线：}$3x^2+4xy+2y^2+8x+4y-6=0$。\hint{使用第 18 题公式，$F_1=3x+2y+4,F_2=2x+2y+2,I_1=5$。}
代入
\[
(3x+2y+4)^2+(2x+2y+2)^2
=5(3x^2+4xy+2y^2+8x+4y-6).
\]
展开整理后得到
\[\boxed{(x+2)^2+(y-1)^2=30.}\]
\ans{轨迹为圆，圆心 $(-2,1)$，半径 $\sqrt{30}$；书后一致。}

\Q{20}{椭圆上的正交切线交点圆（蒙日圆）}
\textbf{椭圆：}$x^2/a^2+y^2/b^2=1$。\hint{记 $F_1=x/a^2,F_2=y/b^2,I_1=1/a^2+1/b^2$，代第 18 题。}
\[
\frac{x^2}{a^4}+\frac{y^2}{b^4}
=\left(\frac1{a^2}+\frac1{b^2}\right)
\left(\frac{x^2}{a^2}+\frac{y^2}{b^2}-1\right).
\]
移项合并后，两侧乘 $a^2b^2$，有
\[\boxed{x^2+y^2=a^2+b^2.}\]
\ans{轨迹是\textbf{与椭圆同心的圆}，半径 $\sqrt{a^2+b^2}$。}

\Q{21}{抛物线正交切线交点轨迹是准线}
\textbf{标准形：}$y^2=2px$（$p>0$），焦点 $(p/2,0)$，准线 $x=-p/2$。\hint{令 $F=y^2-2px$，课本半梯度 $F_1=-p,F_2=y$，迹 $I_1=1$。}
从第 18 题公式得 $p^2+y^2=y^2-2px$，即
\[\boxed{x=-p/2}.\]
因此从准线上的点引出的两条实切线互相垂直；反过来，任意外点正交切线交点必落在准线上。
\ans{准线方程正是正交切线交点轨迹。}

\Q{22}{双曲线正交切线交点轨迹及必须去掉的点}
\textbf{双曲线：}$x^2/a^2-y^2/b^2=1$。\hint{先套第 18 题，之后检查每个候选点引出的切线是否有\textbf{有限实切点}。}
\[
\frac{x^2}{a^4}+\frac{y^2}{b^4}
=\left(\frac1{a^2}-\frac1{b^2}\right)
\left(\frac{x^2}{a^2}-\frac{y^2}{b^2}-1\right)
\Longrightarrow \boxed{x^2+y^2=a^2-b^2}.
\]
当 $a>b$ 时，候选点构成以原点为圆心、半径 $\sqrt{a^2-b^2}$ 的圆。该圆与双曲线的两条渐近线 $y=\pm(b/a)x$ 有\textbf{四个交点}，这些点沿渐近方向只对应无穷远接触，不能形成两条有限实切线，应删除。$a\le b$ 时没有满足题意的外点：$a<b$ 无实圆；$a=b$ 只剩中心，但中心的两条渐近线并非有限切线。
\begin{center}
\begin{tikzpicture}[scale=.88,>=Stealth]
\draw[->](-3.2,0)--(3.3,0)node[right]{$x$};
\draw[->](0,-2.0)--(0,2.0)node[above]{$y$};
\draw[orange!80!black,very thick](0,0)circle(1.73);
\draw[dashed,accent](-3,-1.5)--(3,1.5);
\draw[dashed,accent](-3,1.5)--(3,-1.5);
\draw[thick,mint,domain=-1.0:1.0,samples=70]plot({2*sqrt(1+\x*\x)},\x);
\draw[thick,mint,domain=-1.0:1.0,samples=70]plot({-2*sqrt(1+\x*\x)},\x);
\foreach \xp in {-1.54919,1.54919}{\foreach \yp in {-0.77460,0.77460}{\fill[red!80!black](\xp,\yp)circle(2pt);}}
\node[align=left]at(5.4,.9){示意取 $a=2,b=1$：\\橙色候选圆与两条渐近线\\的四个红色交点必须删去};
\end{tikzpicture}
\end{center}
\ans{$\boxed{x^2+y^2=a^2-b^2}$，在 $a>b$ 时\textbf{去掉与渐近线的四个交点}；其余情况轨迹为空。}

\Q{23}{从椭圆焦点向任一切线作垂足，垂足在辅助圆上}
\textbf{椭圆：}$x^2/a^2+y^2/b^2=1$；焦点 $F_1,F_2$，中心 $O$，$a^2-b^2=c^2$。\hint{使用椭圆光学性质：任一切线关于焦点连线的反射有简洁几何关系。}
设切线在 $P$ 处与椭圆接触，把焦点 $F_1$ 关于切线反射到 $F'_1$，切线的反射性质保证 $F'_1,P,F_2$ 共线、$P$ 在 $F'_1F_2$ 上。故
\[
|F'_1F_2|=|F'_1P|+|PF_2|=|F_1P|+|PF_2|=2a.
\]
若 $H$ 是从 $F_1$ 到切线的垂足，则 $H$ 是 $F_1F'_1$ 的中点；$O$ 又是 $F_1F_2$ 的中点。在三角形 $F_1F'_1F_2$ 中，中位线 $OH$ 长为 $F'_1F_2/2=a$，故
\[\boxed{|OH|=a}.\]
\ans{所有切线垂足都在\textbf{以椭圆中心为圆心、半长轴为半径}的辅助圆上。}

\Q{24}{由一对共轭直径与两点确定圆锥曲线和切线}
\textbf{条件：}共轭直径 $x-y-10=0$ 与 $x+y+6=0$；经过 $(3,-3)$ 和 $(3,-7)$。\hint{两直径垂直且相交于中心 $H=(2,-8)$，在以这两条直径为轴的正交坐标下，无混合项。}
设 $U=x-y-10,\ V=x+y+6$；曲线方程可写为 $pU^2+qV^2=1$。两已知点分别对应 $(U,V)=(-4,6)$ 与 $(0,2)$，得到 $16p+36q=1,\ 4q=1$，即 $q=1/4,p=-1/2$。因此
\[-2U^2+V^2=4.
\]
\readtip{\textbf{等比例方程：}若方程整体乘 $-1$，图像不变。下面采用与书后一致的首项系数为正的写法。}
展开 $-2(x-y-10)^2+(x+y+6)^2-4=0$，整理为
\[\boxed{-x^2+6xy-y^2+52x-28y-168=0}\]
或两边同乘 $-1$：$\boxed{x^2+y^2-6xy-52x+28y+168=0}$。
在点 $P=(3,-3)$ 处 $F_x=2x-6y-52=-28,\ F_y=2y-6x+28=4$，故切线
\[\boxed{7x-y-24=0}.\]
\ans{二次式和切线均与书后最终结果一致；上述系数恢复步骤已展开，避免仅给书后式子。}
